% 华为杯全国研究生数学建模竞赛模板
% 基于 gmcmthesis 文档类
\documentclass[bwprint]{gmcmthesis}

% === 额外宏包（cls 未内置的）===
\usepackage[framemethod=TikZ]{mdframed}
\usepackage[section]{placeins}
\usepackage{needspace}
\usepackage{subfig}
\usepackage{float}
\usepackage{algorithm2e}

% === TikZ ===
\usepackage{tikz}
\usetikzlibrary{arrows.meta,positioning,shapes.geometric,fit,calc}
\usepackage{longtable}

% === 引用上标格式 ===
\setcitestyle{super,square,comma}

% ⛔ 摘要后空白页修复：已迁移到 gmcmthesis.cls（abstract 结尾用 \clearpage
%    替代 \newpage\null，避免内容接近页底时挤出含 \null 的空白页）。
%    此处不需要再 \renewenvironment{abstract}。

% === 竞赛信息 ===
\title{[论文标题]}
\baominghao{[参赛队号]}
\schoolname{[学校名称]}
\membera{[成员A]}
\memberb{[成员B]}
\memberc{[成员C]}

\begin{document}

\maketitle

% === 摘要 ===
\begin{abstract}

[中文摘要内容：问题概述 + 每个子问题的方法和数值结果 + 结论]

\keywords{[关键词1]\quad [关键词2]\quad [关键词3]}
\end{abstract}

\pagestyle{plain}

% === 目录（华为杯篇幅长，建议加目录）===
% ⛔ 用 \clearpage 而非 \newpage：cls 修复后 abstract 已用 \clearpage 换页，
%    这里再用 \clearpage 检测到已在新页开头会自动跳过，不会产生空白页。
\clearpage
\tableofcontents
\clearpage

% === 正文 ===
\input{sections/1_restatement}
\input{sections/2_analysis}
\input{sections/3_assumptions}
\input{sections/4_symbols}
\input{sections/5_problem1}
\input{sections/6_problem2}
\input{sections/7_problem3}
\input{sections/8_sensitivity}
\input{sections/9_evaluation}

% === 参考文献 ===
\begin{thebibliography}{99}
% \bibitem[1]{ref1} 作者. 标题[J]. 期刊, 年份.
\end{thebibliography}

% === 附录 ===
\begin{appendices}
\input{sections/A_code}
\end{appendices}

\end{document}
